\chapter{The Lattice Theory and Reflection Positivity}
\label{ch:lattice}
\label{ch:lattice_reflection_positivity}

\section{The Gauge Group and Configuration Space}
\label{sec:lattice_config}

Let $\Lambda \subset (a \mathbb{Z})^4$ be a four-dimensional hypercubic lattice of lattice spacing $a > 0$, with vertices (sites) $x = (x_0, x_1, x_2, x_3) = (x_0, \mathbf{x}) \in \Lambda$. For finite-volume regularisation, we initially take $\Lambda$ to be a hypercubic torus of spatial length $L = N_s a$ and temporal extent $T = N_\tau a$, with periodic boundary conditions, and eventually consider the thermodynamic limit $\Lambda \nearrow (a \mathbb{Z})^4$.

Let $E(\Lambda)$ denote the set of oriented elementary edges (links) of $\Lambda$. An oriented link $e \in E(\Lambda)$ is uniquely specified by its base vertex $x \in \Lambda$ and a positive coordinate direction $\mu \in \{0, 1, 2, 3\}$, denoted $e = (x, \mu)$. The link connects $x$ to $x + a\hat{\mu}$, where $\hat{\mu}$ is the unit displacement vector in the $\mu$-th direction. The reversed link is denoted $-e = (x + a\hat{\mu}, -\mu)$.

The physical gauge degree of freedom is regularised by assigning to each oriented link $e \in E(\Lambda)$ an element $U_e \in G$, where $G = \SU(N)$ ($N \ge 2$) is the compact, connected, simple special unitary Lie group. Reversing the orientation of an edge corresponds to group inversion:
\begin{equation}
U_{-e} = U_e^{-1} = U_e^\dagger.
\end{equation}
The configuration space of the lattice gauge theory is the compact finite-dimensional Lie group manifold:
\begin{equation}
\Omega_\Lambda = G^{E(\Lambda)} = \prod_{e \in E(\Lambda)} G.
\end{equation}

A local gauge transformation is a map $g : \Lambda \to G$, associating an independent group element $g(x) \in G$ to each lattice site $x \in \Lambda$. The group of all local gauge transformations on $\Lambda$ is the direct product:
\begin{equation}
\mathcal{G}_\Lambda = G^{\Lambda} = \prod_{x \in \Lambda} G.
\end{equation}
The gauge group $\mathcal{G}_\Lambda$ acts smoothly on the configuration manifold $\Omega_\Lambda$ by:
\begin{equation}
\label{eq:lattice_gauge_trans}
(g \cdot U)_e = g(x) \, U_\mu(x) \, g(x + a\hat{\mu})^{-1}, \qquad \text{for } e = (x, \mu).
\end{equation}

\section{Invariant Haar Integration and Peter--Weyl Orthogonality}
\label{sec:haar_measure}

The fundamental probability measure on the unconstrained configuration space $\Omega_\Lambda$ is the product of normalized Haar measures:
\begin{equation}
d\mu_0(U) = \prod_{e \in E(\Lambda)} dU_e.
\end{equation}

The Haar measure $dg$ on the compact Lie group $G = \SU(N)$ is the unique Borel probability measure satisfying:
\begin{enumerate}
    \item \textbf{Left and Right Translation Invariance}: For all $h \in G$ and all continuous test functions $f \in C(G)$:
    \begin{equation}
    \int_G f(h g) \, dg = \int_G f(g h) \, dg = \int_G f(g) \, dg.
    \end{equation}
    \item \textbf{Inversion Invariance}: $\int_G f(g^{-1}) \, dg = \int_G f(g) \, dg$.
    \item \textbf{Normalization}: $\int_G 1 \, dg = 1$.
\end{enumerate}

Because the gauge transformation \eqref{eq:lattice_gauge_trans} acts on each link variable $U_\mu(x)$ by left multiplication by $g(x)$ and right multiplication by $g(x + a\hat{\mu})^{-1}$, the bi-invariance of the Haar measure immediately implies that the product measure $d\mu_0(U)$ is strictly invariant under the entire infinite-dimensional gauge transformation group:
\begin{equation}
d\mu_0(g \cdot U) = d\mu_0(U) \qquad \forall \, g \in \mathcal{G}_\Lambda.
\end{equation}

By the Peter--Weyl theorem \cite{StreaterWightman1964}, the Hilbert space $L^2(G, dg)$ decomposes into an orthogonal direct sum of matrix elements of irreducible representations:
\begin{equation}
L^2(G, dg) = \bigoplus_{R \in \widehat{G}} V_R \otimes V_R^*,
\end{equation}
where $\widehat{G}$ denotes the set of equivalence classes of finite-dimensional irreducible unitary representations of $G$. Let $D^R(g) \in \mathrm{End}(V_R)$ denote the representation matrix in an orthonormal basis, with dimension $d_R = \dim V_R$. The matrix elements satisfy the exact orthogonality relations:
\begin{equation}
\int_G D^R_{ij}(g) \, \overline{D^{R'}_{kl}(g)} \, dg = \frac{1}{d_R} \, \delta_{RR'} \, \delta_{ik} \, \delta_{jl}.
\end{equation}
Taking the trace $i = j$ and $k = l$ yields the character orthogonality theorem:
\begin{equation}
\label{eq:character_orthogonality}
\int_G \chi_R(g) \, \overline{\chi_{R'}(g)} \, dg = \delta_{RR'}, \qquad \chi_R(g) = \Tr_R(D^R(g)).
\end{equation}
This orthogonality guarantees that all character expansions on the lattice are exact and free of metric distortion.

\section{The Wilson Plaquette Action}
\label{sec:wilson_action}

Let $P(\Lambda)$ denote the collection of oriented elementary two-dimensional squares (plaquettes) of the lattice. An oriented plaquette $p = (x, \mu, \nu)$ with $\mu < \nu$ consists of the closed sequence of four directed links:
\begin{equation}
\partial p = (x, \mu) \cup (x + a\hat{\mu}, \nu) \cup (x + a\hat{\nu}, \mu)^{-1} \cup (x, \nu)^{-1}.
\end{equation}
The parallel transport holonomy around the plaquette $p$ is the group element:
\begin{equation}
U_p = U_{\mu\nu}(x) = U_\mu(x) \, U_\nu(x + a\hat{\mu}) \, U_\mu(x + a\hat{\nu})^\dagger \, U_\nu(x)^\dagger \in G.
\end{equation}
Under a local gauge transformation $g \in \mathcal{G}_\Lambda$, the plaquette holonomy transforms by local conjugation at its base vertex $x$:
\begin{equation}
U_p \;\mapsto\; g(x) \, U_p \, g(x)^{-1}.
\end{equation}
Because the trace of a group representation is class-invariant ($\Tr(g A g^{-1}) = \Tr(A)$), the real part of the trace is strictly gauge-invariant:
\begin{equation}
\RePart \, \Tr(U_p) \;\mapsto\; \RePart \, \Tr\bigl( g(x) U_p g(x)^{-1} \bigr) = \RePart \, \Tr(U_p).
\end{equation}

\begin{definition}[Wilson Lattice Action]
The Wilson action $S_W[U]$ on the lattice $\Lambda$ is defined by:
\begin{equation}
\label{eq:wilson_lattice_action}
S_W[U] = \beta \sum_{p \in P(\Lambda)} \left( 1 - \frac{1}{N} \RePart \, \Tr(U_p) \right), \qquad \beta = \frac{2N}{g_0^2},
\end{equation}
where $g_0 > 0$ is the bare gauge coupling constant, and the sum runs over all unoriented elementary plaquettes of $\Lambda$.
\end{definition}

The interacting Euclidean lattice gauge probability measure is:
\begin{equation}
\label{eq:lattice_measure}
d\mu_{\Lambda, \beta}(U) = \frac{1}{\mathcal{Z}_\Lambda(\beta)} \exp\left( -S_W[U] \right) \prod_{e \in E(\Lambda)} dU_e,
\end{equation}
where the partition function is the normalization constant:
\begin{equation}
\mathcal{Z}_\Lambda(\beta) = \int_{\Omega_\Lambda} \exp\left( -S_W[U] \right) \prod_{e \in E(\Lambda)} dU_e.
\end{equation}

\begin{proposition}[Classical Continuum Limit of the Lattice Action]
Let $A_\mu(x)$ be a smooth, classical $\LieSU(N)$-valued connection on $\mathbb{R}^4$, and let link variables be defined by parallel transport:
\begin{equation}
U_\mu(x) = \mathcal{P} \exp\left( i g_0 \int_x^{x + a\hat{\mu}} A_\mu(y) \, dy \right).
\end{equation}
Then as $a \to 0$, the Wilson lattice action reproduces the classical Yang--Mills continuum action with quadratic lattice cutoff errors:
\begin{equation}
S_W[U] = \frac{1}{4} \int_{\mathbb{R}^4} F_{\mu\nu}^a F^{a\mu\nu} \, d^4x + \mathcal{O}(a^2).
\end{equation}
\end{proposition}

\begin{proof}
Applying the Baker--Campbell--Hausdorff formula to the product of four link variables around the square $x \to x + a\hat{\mu} \to x + a\hat{\mu} + a\hat{\nu} \to x + a\hat{\nu} \to x$:
\begin{equation}
U_{\mu\nu}(x) = \exp\left( i g_0 a^2 F_{\mu\nu}(x) + \mathcal{O}(a^3) \right),
\end{equation}
where $F_{\mu\nu} = \partial_\mu A_\nu - \partial_\nu A_\mu + i g_0 [A_\mu, A_\nu]$. Expanding the exponential series:
\begin{equation}
U_{\mu\nu}(x) = I_N + i g_0 a^2 F_{\mu\nu}(x) - \half g_0^2 a^4 F_{\mu\nu}(x)^2 + \mathcal{O}(a^5).
\end{equation}
Taking the normalized trace:
\begin{equation}
\frac{1}{N} \RePart \, \Tr(U_{\mu\nu}(x)) = 1 - \frac{g_0^2 a^4}{2N} \Tr(F_{\mu\nu}(x)^2) + \mathcal{O}(a^6).
\end{equation}
Substituting into \eqref{eq:wilson_lattice_action} and using $\Tr(F_{\mu\nu}^2) = -\half F_{\mu\nu}^a F^{a\mu\nu}$ and $\beta = 2N/g_0^2$:
\begin{align}
S_W[U] &= \frac{2N}{g_0^2} \sum_{x \in \Lambda} \sum_{\mu < \nu} \left( \frac{g_0^2 a^4}{4N} F_{\mu\nu}^a(x) F^{a\mu\nu}(x) + \mathcal{O}(a^6) \right) \nonumber \\
&= \half \sum_{x} a^4 \sum_{\mu < \nu} F_{\mu\nu}^a(x) F^{a\mu\nu}(x) + \mathcal{O}(a^2) \nonumber \\
&= \frac{1}{4} \int_{\mathbb{R}^4} F_{\mu\nu}^a F^{a\mu\nu} \, d^4x + \mathcal{O}(a^2). \qedhere
\end{align}
\end{proof}

\section{Reflection Positivity on the Lattice}
\label{sec:reflection_positivity}

Reflection positivity is the crucial structural link between the Euclidean probabilistic functional measure and the unitarity of relativistic quantum mechanics on complex Hilbert space.

Let $\Pi$ denote a reflection hyperplane in Euclidean spacetime $\mathbb{R}^4$. We distinguish two geometric types of reflection planes:
\begin{enumerate}
    \item \textbf{Site Planes}: Hyperplanes passing through lattice sites, $\Pi_{\mathrm{site}} = \{x \in \mathbb{R}^4 : x_0 = 0\}$.
    \item \textbf{Link Planes}: Hyperplanes passing midway between lattice sites, bisecting temporal links, $\Pi_{\mathrm{link}} = \{x \in \mathbb{R}^4 : x_0 = a/2\}$.
\end{enumerate}
For the construction of the transfer matrix and the physical Hilbert space, we consider reflections across the link plane $\Pi = \{x : x_0 = a/2\}$.

The reflection map $\Theta : \Lambda \to \Lambda$ is defined on vertices by:
\begin{equation}
\Theta(x_0, \mathbf{x}) = (a - x_0, \mathbf{x}).
\end{equation}
The hyperplane divides the lattice into two halves:
\begin{equation}
\Lambda_+ = \left\{ (x_0, \mathbf{x}) \in \Lambda : x_0 \ge a \right\}, \qquad \Lambda_- = \left\{ (x_0, \mathbf{x}) \in \Lambda : x_0 \le 0 \right\} = \Theta(\Lambda_+).
\end{equation}
The set of links $E(\Lambda)$ decomposes into three disjoint subsets:
\begin{equation}
E(\Lambda) = E_+ \cup E_- \cup E_\times,
\end{equation}
where $E_+$ are links contained in $\Lambda_+$, $E_- = \Theta(E_+)$ are links in $\Lambda_-$, and $E_\times = \{( (0, \mathbf{x}), 0 ) : \mathbf{x} \in \mathbb{Z}^3\}$ are temporal links straddling the plane $\Pi$.

On link variables, the reflection operator $\Theta$ acts by:
\begin{equation}
\Theta(U_{(x, \mu)}) = U_{(\Theta x, \mu)}^{-1} = U_{(\Theta x, \mu)}^\dagger \quad \text{if } \mu = 0, \qquad \Theta(U_{(x, j)}) = U_{(\Theta x, j)} \quad \text{if } j \in \{1, 2, 3\}.
\end{equation}
On complex-valued functions $F : \Omega_\Lambda \to \mathbb{C}$, $\Theta$ induces an antilinear involution:
\begin{equation}
(\Theta F)(U) = \overline{F(\Theta U)}.
\end{equation}
Let $\mathcal{A}_+$ denote the commutative algebra of functions depending exclusively on link variables in $E_+ \cup E_\times$.

\begin{theorem}[Osterwalder--Seiler Reflection Positivity \cite{OsterwalderSeiler1978}]
\label{thm:os_rp}
The Wilson lattice gauge measure $d\mu_{\Lambda, \beta}$ is reflection positive with respect to $\Theta$: for any observable $F \in \mathcal{A}_+$:
\begin{equation}
\label{eq:lattice_rp_inequality}
\langle \Theta F, F \rangle_{\mu} = \int_{\Omega_\Lambda} (\Theta F)(U) \, F(U) \, d\mu_{\Lambda, \beta}(U) \ge 0.
\end{equation}
\end{theorem}

\begin{proof}
The total Wilson action $S_W[U]$ decomposes into three parts according to the spatial localization of plaquettes:
\begin{equation}
S_W[U] = S_+[U] + S_-[U] + S_\times[U],
\end{equation}
where $S_+[U]$ contains plaquettes entirely in $\Lambda_+$, $S_-[U] = \Theta S_+[U]$ contains plaquettes in $\Lambda_-$, and $S_\times[U]$ contains temporal plaquettes intersecting the reflection plane $\Pi$. 

A temporal plaquette intersecting $\Pi$ connects site $(0, \mathbf{x})$ to $(a, \mathbf{x})$ and $(0, \mathbf{x} + a\hat{j})$ to $(a, \mathbf{x} + a\hat{j})$. In temporal gauge on the straddling links ($U_0(0, \mathbf{x}) = I_N$), the plaquette holonomy simplifies to:
\begin{equation}
U_p = U_j(0, \mathbf{x}) \, U_j(a, \mathbf{x})^\dagger = U_j(0, \mathbf{x}) \, (\Theta U_j(0, \mathbf{x}))^\dagger.
\end{equation}
Expanding the Boltzmann weight of the straddling interaction using the Peter--Weyl character expansion:
\begin{equation}
\exp\left( -S_\times[U] \right) = \prod_{\mathbf{x}, j} \sum_{R \in \widehat{G}} c_R(\beta) \sum_{k, l=1}^{d_R} D^R_{kl}\bigl( U_j(0, \mathbf{x}) \bigr) \, \overline{D^R_{kl}\bigl( U_j(a, \mathbf{x}) \bigr)}.
\end{equation}
Because $c_R(\beta) > 0$ for all representations $R$, this factorizes into a sum of positive terms of the form $\sum_i \overline{A_i(\Theta U)} A_i(U)$:
\begin{equation}
\exp\left( -S_\times[U] \right) = \sum_\alpha \Theta(H_\alpha)[U] \, H_\alpha[U],
\end{equation}
where each $H_\alpha$ depends only on links in $\Lambda_+ \cup E_\times$. Therefore, for any $F \in \mathcal{A}_+$:
\begin{align}
\int (\Theta F) F \, e^{-S_W} \, d\mu_0 &= \sum_\alpha \int (\Theta F) F \, \Theta(e^{-S_+}) e^{-S_+} \, \Theta(H_\alpha) H_\alpha \, d\mu_0 \nonumber \\
&= \sum_\alpha \int_{G^{E_\times}} dU_\times \left| \int_{G^{E_+}} F \, e^{-S_+} \, H_\alpha \, \prod_{e \in E_+} dU_e \right|^2 \ge 0. \qedhere
\end{align}
\end{proof}

\section{The Physical Hilbert Space and the Transfer Matrix}
\label{sec:hilbert_transfer}

Let $\mathcal{N}_a = \{F \in \mathcal{A}_+ : \langle \Theta F, F \rangle_\mu = 0\}$ denote the null space of the reflection-positive semi-definite form. The pre-Hilbert space is the quotient:
\begin{equation}
\mathcal{D}_a = \mathcal{A}_+ / \mathcal{N}_a,
\end{equation}
and the physical Hilbert space $\mathcal{H}_a$ at lattice spacing $a$ is its metric completion:
\begin{equation}
\mathcal{H}_a = \overline{\mathcal{A}_+ / \mathcal{N}_a}.
\end{equation}
The inner product on $\mathcal{H}_a$ is given by $\langle [F], [G] \rangle_{\mathcal{H}} = \langle \Theta F, G \rangle_\mu$.

The transfer matrix operator $\hat{T}_a : \mathcal{H}_a \to \mathcal{H}_a$ is the operator advancing states by one temporal lattice unit $a$. As established in Theorem \ref{thm:os_rp}, $\hat{T}_a$ is a bounded, positive, self-adjoint contraction on $\mathcal{H}_a$:
\begin{equation}
0 \le \hat{T}_a \le \mathbf{1}, \qquad \hat{T}_a^\dagger = \hat{T}_a.
\end{equation}
The discrete-lattice Hamiltonian $\hat{H}_a$ is defined by:
\begin{equation}
\hat{T}_a = \exp(-a \hat{H}_a) \iff \hat{H}_a = -\frac{1}{a} \ln \hat{T}_a \ge 0.
\end{equation}
This completes the non-perturbative construction of the quantum state space and time-evolution operator at finite lattice spacing.
